\section{四条路线在最小化哪个散度}
\label{sec:14-Deep-Learning/07-Generative-Models-Unified/01}

组里两个人各交了一个人脸生成模型。第一个是 VAE，生成的图全是糊的，五官像隔了一层毛玻璃，但连着采两百张，肤色、年龄、发型、戴不戴眼镜都出现过。第二个是 GAN，随便挑一张都锐利到可以直接放进 PPT，可采到第三十张你就开始眼熟——翻来覆去是四五张脸换了角度。

会上的解释是训练不够：VAE 再多跑几个 epoch 就清晰了，GAN 换个学习率就多样了。两条都跑过，两条都没用。原因不在超参数。这两种失败分别写在两个训练目标里，而且方向恰好相反——一个被目标函数逼着去覆盖所有它见过的东西，哪怕代价是在不该有东西的地方也放上质量；另一个被目标函数允许对没生成出来的模式完全不负责，只要生成出来的那些足够像。

这一篇要把这句话变成可以推的式子。式子只有一个不对称，后面四条路线的性格都从它长出来。

\subsection{一个写得出但算不了的目标}

生成建模想做的事写下来只有一行：给定真实数据分布 \(p_{\text{data}}\)，在一个参数族里挑出最接近它的 \(p_\theta\)，

\[
\min_\theta\ D(p_{\text{data}}\,\Vert\,p_\theta).
\]

麻烦是这一行里没有一个符号是可以直接用的。\(p_{\text{data}}\) 你手上只有从它抽出的样本，密度值一个也没有；\(p_\theta\) 这一侧更糟——大量有表达力的模型能写出一个未归一化的打分 \(\tilde p_\theta(x)\)，却算不出归一化常数

\[
Z_\theta=\int \tilde p_\theta(x)\,dx.
\]

高维空间上的这个积分和 \hyperref[sec:14-Deep-Learning/06-Variational-Autoencoders/01]{``VAE 把不可积后验变成可训练下界''} 里挡住边缘似然的那个积分是同一类困难，详细的账在 \hyperref[sec:13-Machine-Learning/08-Sampling-and-Inference/06]{``配分函数把能量变成全局概率''} 里算过。

所以真正要做的选择有两层：先选 \(D\)，再选一条绕开不可算量的路。第二层是技术，第一层是立场——它决定模型在数据分布覆盖不到的地方怎么办，而生成模型的成败恰恰全在那些地方。

\subsection{同一个 KL，两个方向，两种失败}

先把 \(D\) 取成 KL 散度。它不对称，于是有两个写法，两个写法给出的最优解可以差得面目全非。

前一个方向，数据在外：

\[
\KL(p_{\text{data}}\,\Vert\,p_\theta)
=\int p_{\text{data}}(x)\,\log\frac{p_{\text{data}}(x)}{p_\theta(x)}\,dx.
\]

看被积函数在哪里会爆炸。权重是 \(p_{\text{data}}(x)\)。凡是 \(p_{\text{data}}(x)>0\) 的地方，只要 \(p_\theta(x)\to0\)，\(\log\) 项就趋于 \(+\infty\)，而它前面挂着一个非零权重——整个积分发散。翻译成一句话：\textbf{模型敢在有数据的地方留空洞，代价是无穷大}。反过来，在 \(p_{\text{data}}(x)=0\) 的地方，无论 \(p_\theta(x)\) 堆了多少质量，被积函数都是 \(0\)，一分钱不罚。

后一个方向，模型在外：

\[
\KL(p_\theta\,\Vert\,p_{\text{data}})
=\int p_\theta(x)\,\log\frac{p_\theta(x)}{p_{\text{data}}(x)}\,dx.
\]

权重换成了 \(p_\theta(x)\)，惩罚的位置整个翻过来。凡是 \(p_\theta(x)>0\) 而 \(p_{\text{data}}(x)\to0\) 的地方，积分发散——\textbf{模型敢在没数据的地方放质量，代价是无穷大}。而某个真实模式上 \(p_\theta\) 直接取零？那里被积函数是 \(0\cdot\log 0\)，按约定等于 \(0\)，同样一分钱不罚。

两句话摆在一起就是全部：前向 KL 罚漏掉，不罚多出来；反向 KL 罚多出来，不罚漏掉。这不是两个近似程度不同的量，是两套完全不同的赏罚制度。KL 的方向性本身在 \hyperref[sec:07-Information-Theory/03-Divergences/02]{``Gibbs 不等式与 KL 的方向性''} 里已经从编码代价推过一遍，这里要用的是它在参数族被迫做取舍时的后果。

取舍在参数族容不下真分布时才真正暴露。设数据是两个分开的窄峰，而你只允许用一个高斯去拟合。

\begin{example}
\(p_{\text{data}}=\tfrac12\mathcal N(-a,\sigma_0^2)+\tfrac12\mathcal N(a,\sigma_0^2)\)，模型族是全体一维高斯 \(\mathcal N(\mu,\sigma^2)\)。

最小化 \(\KL(p_{\text{data}}\Vert p_\theta)\) 等价于最大化 \(\E_{p_{\text{data}}}[\log p_\theta]\)，对高斯而言这就是矩匹配：最优解取 \(\mu=0\)、\(\sigma^2=a^2+\sigma_0^2\)。当 \(a\) 比 \(\sigma_0\) 大得多时，这个高斯的顶点正落在两峰之间的谷底——数据在那里几乎没有样本，模型却把最高的密度放在了那里。

最小化 \(\KL(p_\theta\Vert p_{\text{data}})\) 则得到两个对称的局部极小，分别贴在 \(\mu\approx\pm a\) 上，\(\sigma\approx\sigma_0\)。梯度下降落到哪一个取决于初始化，落定之后另一个峰彻底不在模型的支撑里。
\end{example}

\begin{center}
\begin{tikzpicture}[x=1.42cm,y=1.42cm,>=stealth]
  \draw[->] (-3.4,0) -- (3.6,0) node[right] {\(x\)};
  \draw[->] (-3.4,0) -- (-3.4,2.5) node[above] {密度};
  % 数据分布：双峰
  \draw[black,thick] plot[domain=-3.3:3.3,samples=140]
    (\x,{exp(max(-(\x+1.5)*(\x+1.5)/0.605,-9))+exp(max(-(\x-1.5)*(\x-1.5)/0.605,-9))});
  % 前向 KL：摊平覆盖
  \draw[blue!70!black,thick,dashed] plot[domain=-3.3:3.3,samples=140]
    (\x,{0.78*exp(max(-\x*\x/3.645,-9))});
  % 反向 KL：贴住一峰
  \draw[red!75!black,thick,dash dot] plot[domain=-3.3:3.3,samples=140]
    (\x,{2.0*exp(max(-(\x-1.5)*(\x-1.5)/0.605,-9))});
  % 谷底的多余质量
  \draw[<->,gray!85] (0,0.03) -- (0,0.75);
  \node[below,font=\small,align=center] at (0,-0.1)
    {谷底没有数据，\\前向 KL 的解仍在这里放质量};
  % 图例
  \draw[black,thick] (-3.15,2.28) -- (-2.65,2.28);
  \node[right,font=\small] at (-2.6,2.28) {\(p_{\text{data}}\)：两个分开的模式};
  \draw[blue!70!black,thick,dashed] (-3.15,1.92) -- (-2.65,1.92);
  \node[right,font=\small] at (-2.6,1.92) {前向 KL 的解：摊开覆盖两峰};
  \draw[red!75!black,thick,dash dot] (-3.15,1.56) -- (-2.65,1.56);
  \node[right,font=\small] at (-2.6,1.56) {反向 KL 的解：只贴住一峰};
\end{tikzpicture}

\emph{同一份双峰数据，两个 KL 方向选出两个几乎不相干的解：一条摊开去够住两个峰、代价是在谷底留下本不该有的质量；一条只贴住一个峰、代价是另一个峰整个消失。}
\end{center}

这张图后面几篇会反复用到。现在把它翻译回开头那两个模型。

蓝色那条曲线上的样本长什么样？在两峰之间采到的点既不像左峰也不像右峰，是两类东西的中间态。图像上，这就是模糊：模型不敢在任何一种清晰形态上把质量收紧，因为收紧意味着别的形态附近出现空洞，而空洞在前向 KL 下罚得倾家荡产。于是它输出所有可能形态的某种加权模糊平均。红色那条曲线上的样本每一张都锐利，因为它们全部来自一个被贴紧的真实模式——代价是另外那个峰上的数据在模型眼里根本不存在。这就是模式坍缩。

\begin{remark}
模糊和模式坍缩常被当作两种工程故障，各自配一套调参偏方。它们其实是同一个不对称的两侧：一个是覆盖式目标在参数族不够大时的必然产物，一个是模式寻求式目标在允许放弃的前提下的最优解。不加约束地``修好''其中一个，通常只是把模型推到另一侧去。
\end{remark}

需要说清楚的边界：上面的对照建立在参数族容不下 \(p_{\text{data}}\) 这个前提上。如果模型族足够大、优化足够充分、数据足够多，两个方向的最优解都是 \(p_{\text{data}}\) 本身，差别消失。现实中三个条件没有一个成立，所以差别一直在。

\subsection{四条路线各站在哪个位置}

有了这把尺子，四条路线可以一次归位。

\textbf{最大似然是前向 KL。} 把 \(p_{\text{data}}\) 换成 \(n\) 个样本的经验分布，

\[
\begin{aligned}
\argmin_\theta\ \KL(p_{\text{data}}\,\Vert\,p_\theta)
&=\argmax_\theta\ \E_{p_{\text{data}}}\!\left[\log p_\theta(x)\right]\\
&\approx\argmax_\theta\ \frac1n\sum_{i=1}^n\log p_\theta(x_i),
\end{aligned}
\]

因为 \(\E_{p_{\text{data}}}[\log p_{\text{data}}]\) 与 \(\theta\) 无关，可以从目标里扔掉。这条恒等在 \hyperref[sec:13-Machine-Learning/11-Information-Theoretic-Learning/01]{``交叉熵训练其实在最小化 KL 差异''} 里已经建立。凡是训练目标写成对数似然的模型，都自动继承覆盖式的性格。

\textbf{VAE 优化的是似然的下界。} ELBO 与对数似然之间差一个间隙，

\[
\log p_\theta(x)=\mathcal L(\theta,\phi;x)+\KL\!\left(q_\phi(z\given x)\,\Vert\,p_\theta(z\given x)\right).
\]

主目标在前向 KL 的方向上，所以 VAE 在数据层面是覆盖式的——这解释了模糊。有意思的是间隙那一项：它是近似后验对真实后验的\emph{反向}KL，于是在潜空间那一层，变分推断反而是模式寻求的（\hyperref[sec:13-Machine-Learning/08-Sampling-and-Inference/03]{``变分推断把后验近似变成优化''} 讲的正是这一层）。同一个模型的两层用了相反方向的 KL，一层负责别漏掉数据，一层负责别把后验估宽。

\textbf{归一化流直接优化精确似然。} 它同样站在前向 KL 的方向上，但没有间隙——通过可逆映射和变量替换公式，\(\log p_\theta(x)\) 是闭式可算的。代价换到了架构上：映射必须双射、维数不能变、Jacobian 行列式必须便宜。下一篇处理这笔交易。

\textbf{GAN 的原始目标在最优判别器下变成 JS 散度。} 训练完的判别器不是对手，是散度的估计器；把最优判别器代回目标得到

\[
2\,\operatorname{JS}(p_{\text{data}}\,\Vert\,p_\theta)-\log 4 .
\]

JS 是对称的，两个方向的惩罚都在（\hyperref[sec:07-Information-Theory/03-Divergences/03]{``Jensen--Shannon 散度与尺度选择''}）。对称并不意味着中庸：JS 在两个分布支撑几乎不重叠时会饱和成常数，梯度消失，实际训练中生成器往往退到反向 KL 那一侧的行为。第 03 篇把这条链推完。

\textbf{扩散模型训练的是一个加权去噪回归。} 表面上看它既不是似然也不是判别器博弈，只是让网络从加噪样本里预测噪声。但这个平方误差目标可以被证明等于一个变分下界的加权形式，因而仍在前向 KL 的方向上。它的特殊之处是把一次跨越整个分布的匹配，沿噪声尺度切成了一串条件分布的匹配，每一次匹配的两个分布都近得多。

四条路线的目标函数长得毫不相似，落在这把尺子上却只有三个位置：三条在前向 KL 一侧（区别是精确、下界、还是分解成多步），一条在对称的 JS 上（区别是散度本身要靠一个网络估出来）。开头那两个模型的病症，到这里已经不需要额外解释。

\subsection{归一化常数：四种绕法}

还有一个共同的技术障碍需要点破，它和方向的选择正交。

能量模型把话说得最直白：给每个 \(x\) 一个能量 \(E_\theta(x)\)，密度是 \(p_\theta(x)=e^{-E_\theta(x)}/Z_\theta\)。这个写法极其自由，网络想多复杂都行。求对数似然的梯度：

\[
\nabla_\theta\log p_\theta(x)
=-\nabla_\theta E_\theta(x)-\nabla_\theta\log Z_\theta
=-\nabla_\theta E_\theta(x)+\E_{p_\theta}\!\left[\nabla_\theta E_\theta(x')\right].
\]

第一项在数据点上算，便宜。第二项是在\emph{模型自己的分布}下取期望——想估它就得从 \(p_\theta\) 采样，而从一个只知道未归一化密度的分布采样正是 MCMC 要解决的难题，每一步训练都要跑一次。自由能与配分函数这套账在 \hyperref[sec:13-Machine-Learning/13-Statistical-Physics-and-Free-Energy/01]{``自由能把配分函数变成热力学量''} 里有完整交代。

四条路线各有各的绕法，而且四种绕法互不相同：

VAE 不去碰 \(Z_\theta\)，它把难算的量放进一个下界里，用 Jensen 不等式把对数搬到期望内部，从此只需要能采样的近似后验。归一化流让 \(Z_\theta\) 恒等于 \(1\)：基分布已经归一化，可逆映射保持总质量，变量替换公式把密度精确搬运过去，没有任何常数需要估计。GAN 干脆不建模密度——它只学一个从噪声到样本的确定性映射，密度值从头到尾没出现过，代价是你也就永远拿不到 \(\log p_\theta(x)\)。扩散模型把一次困难的匹配拆成很多次容易的匹配：每一步只需要匹配一个加了少量噪声的条件分布，而这种局部匹配可以化成回归，回归不需要归一化常数。

一句话概括：可算的似然、可采样的模型、任意灵活的架构，这三样东西没有一条路线能同时拿全。选哪条路线，本质上是在决定放弃哪一样。

下一篇从最不肯妥协的那条开始——归一化流坚持要精确似然，看看为这份精确，架构上要交出多少。
